% !TeX root = ../动手学强化学习笔记.tex
% 由主文档通过 \input 载入；请编译主文档。
\setcounter{chapter}{22} % 保留原书章号。
\chapter{多智能体强化学习}
\label{ch:rl-multiagent}

多智能体强化学习（multi-agent reinforcement learning）中，其他决策者也可能学习和改变策略。一个学习者面对的环境由物理转移与其他策略共同决定，因此单智能体平稳性假设和单一最优性概念都需要重新检查。本章以矩阵博弈和随机博弈建立基本语言，区分团队收益、最佳响应与均衡，再介绍独立学习、集中训练与分散执行、自博弈的机制及评价要求。学习重点是辨认决策者之间的相互影响和训练、执行的信息边界，通用均衡求解及共同学习动态的收敛理论留待进一步学习。

\section{矩阵博弈与随机博弈}

有限矩阵博弈给出各参与者的动作集和对每个联合动作的收益。双人零和博弈中，行玩家收益矩阵为 $M$，列玩家收益为其负值；若混合策略分别为 $p,q$，期望收益为 $p^\top Mq$。最佳响应是在对手策略固定时最大化自身收益的策略。纳什均衡（Nash equilibrium）是一组策略，使任何参与者单独改变自身策略都不能提高其收益。

\begin{example}[为什么均衡可能需要随机化]
取 $M=\left(\begin{smallmatrix}1&-1\\-1&1\end{smallmatrix}\right)$。若行玩家总选第一行，列玩家可选第二列使其收益为 $-1$；任何纯策略都可被利用。双方各以 $1/2$ 选择两个动作时，任一方单独改变概率都不能提高收益，故构成收益为零的均衡。若第一动作概率为 $p,q$，行收益为 $(2p-1)(2q-1)$，可直接验证该性质。
\end{example}

随机博弈（stochastic game）进一步加入状态 $s$、联合动作 $\boldsymbol a$、转移 $P(s'\mid s,\boldsymbol a)$ 和每个智能体自己的奖励 $r_i$。合作问题共享团队收益，一般和问题则允许利益部分一致、部分冲突；团队最优、单方最佳响应和均衡是不同评价标准。

\section{独立学习与非平稳性}

若其他智能体策略固定，当前智能体可以把它们的动作平均进环境转移。以两个智能体为例，诱导转移为 $P^{\pi_2}(s'\mid s,a_1)=\sum_{a_2}\pi_2(a_2\mid s)P(s'\mid s,a_1,a_2)$。当 $\pi_2$ 在训练中改变，这个转移核也改变，过去回放数据可能不再代表当前对手。

独立 Q-learning 把其他智能体视为环境的一部分，实现简单，但同时更新时可能出现追逐、振荡和失效的价值估计。对手冻结、交替更新或保存历史对手可缓和某些问题，但不能只凭单智能体收敛定理证明整个共同学习过程收敛。

\section{集中训练与分散执行}

集中训练、分散执行（centralized training with decentralized execution，CTDE）允许训练时使用额外联合信息构造评论家，执行时各演员仍仅使用自己的局部观测历史。例如评论家可输入全局状态和联合动作，演员 $\pi_i(a_i\mid h_i)$ 则只能读取本地历史 $h_i$。

这种结构帮助评论家区分其他参与者的影响，但训练阶段允许使用的信息必须事先规定。共享团队奖励还带来信用分配问题：团队成功不能直接说明每个局部动作都有效。反事实基线可以比较“保持其他条件不变时本智能体不同动作”的评价，但估计依赖评论家是否准确，并非真实实验中自动可观测的反事实。

\section{自博弈与评价对象}

自博弈（self-play）把当前或历史策略作为对手生成训练数据。只与某个固定版本比赛，可能学会专门利用其弱点；与多版本对手训练和交叉对战有助于暴露这种过拟合。零和问题还可考察对最佳响应的脆弱程度，但大环境中最佳响应通常只能近似求得，评价需说明这一限制。

合作问题要区分固定队友配合与新队友适应。跨不同团队组合的成绩比单一训练搭档更能揭示协调规则是否可迁移。对某个对手高胜率、平均对战较好和接近均衡不是可以互换的结论。

\section{习题与实验}
\begin{enumerate}
\item 验证矩阵例子的均衡；若对手选择第一动作的概率为 $0.7$，求行玩家的最佳响应。
\item 构造双方依次采用纯最佳响应而持续循环的过程，解释最佳响应合理为何不意味着共同动态收敛。
\item 在合作网格中比较独立评论家和集中评论家，保持执行信息一致，并在新队友组合上评价。训练与测试所用对手集合应分别列明。
\end{enumerate}





